\documentclass[12pt]{article}
\usepackage[spanish]{babel}
\usepackage{amsmath,amssymb}
\usepackage{geometry}
\geometry{margin=2.5cm}

\title{Integrales Impropias}
\author{}
\date{}

\begin{document}
\maketitle

\section*{Ejercicio 38}
Calcular la integral impropia:
\[
\int_{1}^{+\infty} \frac{dx}{x^2}
\]

La integral es impropia debido a que el límite superior es infinito. Se procede a evaluarla como un límite.

\[
\int \frac{dx}{x^2} = \int x^{-2} dx = -x^{-1}
\]

Por lo tanto:
\[
\int_{1}^{+\infty} \frac{dx}{x^2}
= \lim_{b \to +\infty} \left[ -\frac{1}{x} \right]_{1}^{b}
\]

\[
= \lim_{b \to +\infty} \left(-\frac{1}{b} + 1\right) = 1
\]

La integral es convergente y su valor es:
\[
\boxed{1}
\]

\section*{Ejercicio 39}
Calcular la integral impropia:
\[
\int_{-\infty}^{0} \frac{dx}{(4-x)^2}
\]

La integral es impropia debido al límite inferior infinito. Se realiza la sustitución:
\[
u = 4 - x \Rightarrow du = -dx
\]

Entonces:
\[
\int \frac{dx}{(4-x)^2} = -\int u^{-2} du = \frac{1}{u}
\]

Volviendo a la variable original:
\[
\frac{1}{4-x}
\]

Evaluando:
\[
\int_{-\infty}^{0} \frac{dx}{(4-x)^2}
= \lim_{a \to -\infty} \left[ \frac{1}{4-x} \right]_{a}^{0}
\]

\[
= \frac{1}{4} - \lim_{a \to -\infty} \frac{1}{4-a}
= \frac{1}{4}
\]

La integral es convergente y su valor es:
\[
\boxed{\frac{1}{4}}
\]

\section*{Ejercicio 40}
Calcular la integral impropia:
\[
\int_{0}^{1} \frac{dx}{\sqrt{x}}
\]

Esta integral es impropia porque la función no está definida en $x = 0$.

\[
\frac{1}{\sqrt{x}} = x^{-1/2}
\]

\[
\int x^{-1/2} dx = 2x^{1/2}
\]

Evaluando mediante límites:
\[
\int_{0}^{1} \frac{dx}{\sqrt{x}}
= \lim_{a \to 0^+} \left[2\sqrt{x}\right]_{a}^{1}
\]

\[
= 2(1 - 0) = 2
\]

La integral es convergente y su valor es:
\[
\boxed{2}
\]

\end{document}